\documentclass[12pt,a4paper]{article}

\usepackage[romanian]{babel}
\usepackage[T1]{fontenc}
\usepackage[utf8]{inputenc}
\usepackage{lmodern}
\usepackage{enumitem}

\title{Modelul Relațional și Regulile lui Codd}
\author{}
\date{}

\begin{document}

\maketitle

\section{Modelul Relațional}
Modelul relațional, introdus de E.F. Codd, reprezintă datele sub forma unor relații
(tabele), fiecare tabel fiind compus din:
\begin{itemize}[noitemsep]
    \item tupluri (rânduri),
    \item atribute (coloane),
    \item o cheie primară,
    \item reguli de integritate.
\end{itemize}

Acest model stă la baza tuturor sistemelor moderne de gestiune a bazelor de date.

\section{Regulile lui Codd (0--12)}
Codd a formulat 13 reguli care definesc cerințele unui SGBD relațional complet.

\begin{enumerate}[label=\textbf{Regula \arabic*:}, leftmargin=2cm]
    \item \textbf{Regula fundamentală (0):} Sistemul trebuie să gestioneze datele exclusiv prin capabilități relaționale.
    \item \textbf{Regula informației:} Toate datele sunt reprezentate sub formă de valori în tabele.
    \item \textbf{Acces garantat:} Fiecare valoare atomică este accesibilă prin tabel--cheie--coloană.
    \item \textbf{Tratamentul valorilor NULL:} NULL trebuie tratat uniform, distinct de zero sau șirul vid.
    \item \textbf{Catalog online activ:} Metadatele trebuie stocate în tabele accesibile prin limbajul relațional.
    \item \textbf{Limbaj de date complet:} Sistemul trebuie să ofere un limbaj complet (ex.: SQL).
    \item \textbf{Actualizarea vizualizărilor:} View-urile actualizabile trebuie să poată fi actualizate.
    \item \textbf{Operații la nivel de set:} Manipularea datelor trebuie să funcționeze pe seturi, nu doar pe elemente individuale.
    \item \textbf{Independență fizică:} Modificările fizice nu afectează aplicațiile.
    \item \textbf{Independență logică:} Modificările logice nu afectează aplicațiile existente.
    \item \textbf{Independența integrității:} Regulile de integritate trebuie definite în baza de date.
    \item \textbf{Independența distribuției:} Sistemul trebuie să funcționeze identic în medii distribuite.
    \item \textbf{Nonsubversion:} Niciun mecanism de nivel inferior nu trebuie să ocolească regulile relaționale.
\end{enumerate}

\end{document}
